\documentclass[14pt,a4paper]{extarticle}
\usepackage[left=3cm,right=2cm,top=2cm,bottom=2cm,headheight=18pt]{geometry}
\usepackage{fontspec}
\setmainfont{Liberation Serif}
\setsansfont{Liberation Serif}
\setmonofont{Liberation Mono}
\usepackage{unicode-math}
\setmathfont{texgyretermes-math.otf}
\usepackage[english,russian]{babel}
\usepackage{amsmath,graphicx,booktabs,longtable,array,enumitem}
\usepackage[normalem]{ulem}
\usepackage{titlesec,fancyhdr}
\usepackage[hidelinks,unicode]{hyperref}
\usepackage{url}
\urlstyle{same}
\setlength{\parindent}{1.27cm}
\setlength{\parskip}{0pt}
\linespread{1}
\setlength{\emergencystretch}{3em}
\clubpenalty=10000
\widowpenalty=10000
\titleformat{\section}{\normalfont\bfseries}{\thesection.}{0.5em}{}
\titlespacing*{\section}{0pt}{1.1em}{0.5em}
\setlist{nosep,leftmargin=1.27cm}
\pagestyle{fancy}
\fancyhf{}
\renewcommand{\headrulewidth}{0pt}
\renewcommand{\footrulewidth}{0pt}
\fancyfoot[C]{\ifnum\value{page}>1\relax\thepage\fi}
% Information-analytical material: O‘RQ-682 appended methodology §83 (11–14 pt), contents §84, comparison table §85.
\renewcommand{\normalsize}{\fontsize{14bp}{16.8bp}\selectfont}
\normalsize
\newcommand{\tsize}{\fontsize{12bp}{14.4bp}\selectfont}
\newcolumntype{L}[1]{>{\raggedright\arraybackslash}p{#1}}
\setlength{\LTcapwidth}{\linewidth}
\begin{document}
\begin{flushright}
Директору Национального агентства\\
социальной защиты при Президенте\\
Республики Узбекистан\\[2mm]
Копия: Министерству экономики и финансов\\
Республики Узбекистан\\[4mm]
от Ашуралиева Абдухолика Акмал угли,\\
независимого исследователя\\
(ORCID 0009-0003-5482-5526)\\
Место жительства: 100017, г. Ташкент,\\
Юнусабадский район, массив Кашгар, дом 16\\
\href{mailto:abdulxoliqashuraliyev@gmail.com}{abdulxoliqashuraliyev@gmail.com}\\[2mm]
8 октября 2026 года
\end{flushright}

\begin{center}\textbf{ПРЕДЛОЖЕНИЕ}\par
\textbf{к проекту нормативно-правового акта, разрабатываемому во исполнение пункта 14 Указа Президента Республики Узбекистан от 11 сентября 2026 года № УП-193}\end{center}

В соответствии со статьёй 40 Конституции Республики Узбекистан, статьёй 20 Закона «О нормативно-правовых актах» и статьями 3 и 5 Закона «Об обращениях физических и юридических лиц» вношу предложения для включения в проект акта о новой системе включения семей в Социальный реестр на основе многомерной оценки, который Агентство совместно с Министерством экономики и финансов вносит в Администрацию Президента до 1 декабря 2026 года.

С 1 января 2027 года категория семьи будет определяться балльной системой; по действующему Положению итоговая оценка и категория семьи формируются информационным модулем автоматически. Предложения обеспечивают, чтобы результат оценки каждой семьи вытекал только из утверждённых правил, сохранялся с основаниями и объяснялся заявителю, как требуют статья 24 Закона «О персональных данных» и статьи 53–54 Закона «Об административных процедурах», а также устраняют несогласованности действующего Положения о порядке ведения Социального реестра, доказанные расчётами. Предлагаемые нормы не создают новых органов и информационных систем; их реализация предлагается в пределах существующих средств при условии оценки Агентством (раздел 10 материала 2). Персональные данные семей не запрашиваются.

Прошу рассмотреть и дать ответ по существу по каждому из следующих вопросов:
\begin{enumerate}[leftmargin=1.27cm,label=\arabic*.]
\item Включить в проект положения раздела I материала 1 (пункты 1–23) либо нормы, обеспечивающие те же результаты.
\item Устранить случай, при котором семья с доходом, \textit{равным} минимальным потребительским расходам, признаётся по пункту 32 Положения, но не относится ни к одной категории пункта 39, — изменением подпункта «а» пункта 33 (раздел II, пункт 3 материала 1).
\item Исключить отрицательную площадь приусадебного участка в формуле пункта 27 Положения (раздел II, пункт 2).
\item Исключить уменьшение учитываемого дохода при увеличении фактического денежного перевода и указать период сравнения в пункте 26 Положения (раздел II, пункт 1), выбрав вариант по результатам оценки совокупного воздействия двух вариантов, указанных в разделе 10 материала 2.
\item Обеспечить предоставление заявителю конкретных фактических и юридических оснований решения о включении, отказе, изменении категории или исключении, включая применённую норму, порядка возражения, срока и порядка обжалования в соответствии со статьёй 24 Закона «О персональных данных» и статьями 53–54 Закона «Об административных процедурах» (раздел II, пункты 4 и 5). Если Агентство считает, что эта обязанность уже исполняется, прошу указать норму и способ, которым заявитель получает эти сведения.
\item Согласовать пункты 34–35 и подпункт «в» пункта 39 Положения с подпунктом «в» пункта 3 Указа № УП-258 и пунктом (i) примечания (b) к приложению № 1 к нему: Положение включает в категорию «семья, находящаяся в черте бедности» семьи в тяжёлой жизненной ситуации, уменьшенный доход которых ниже минимальных потребительских расходов, что Указом не предусмотрено (раздел 7 материала 2).
\item Учесть настоящее предложение при подготовке проекта в соответствии со статьёй 23 Закона «О нормативно-правовых актах» и сообщить мне о размещении проекта на портале обсуждения проектов нормативно-правовых актов в соответствии со статьёй 24 этого Закона, а также указать подразделение Агентства для взаимодействия по представленным расчётам.
\end{enumerate}

Согласно статье 27 Закона «Об обращениях физических и юридических лиц» обращение считается рассмотренным, если рассмотрены все поставленные в нём вопросы, а ответ содержит конкретные основания по каждому вопросу. Поскольку проект вносится до 1 декабря 2026 года, прошу рассмотреть предложение в месячный срок, установленный статьёй 28 этого Закона, а если потребуется дополнительное изучение — уведомить меня об этом в десятидневный срок, как предусмотрено той же статьёй. Прошу сообщить, какие положения включены в проект, в какой редакции, либо какой действующей нормой вопрос уже решён, либо по какой причине положение не принято.

Все расчёты воспроизводятся приложенной программой. Набор контрольных примеров и программа передаются безвозмездно; готов продемонстрировать их ответственным работникам.

\noindent Приложения:
\begin{enumerate}[leftmargin=1.27cm,label=\arabic*.]
\item Материал 1. Предложения для включения в проект нормативно-правового акта.
\item Материал 2. Пояснительная записка с научно-правовым обоснованием.
\item Электронное приложение: программа контрольных расчётов, тесты и результаты.
\end{enumerate}
Основная редакция материалов — русская; узбекская и английская редакции приложены. Узбекская редакция материала 1 содержит формулировки изменений на государственном языке.

\bigskip
\noindent Ашуралиев Абдухолик Акмал угли
\end{document}
